\documentclass[10pt,a4paper]{amsart}
\usepackage[margin=19mm]{geometry}
\usepackage{amsmath,amssymb,mathtools}
\usepackage[scr=rsfs,cal=euler]{mathalfa}
\usepackage{xfrac}
\usepackage[unicode,hidelinks,bookmarks=false]{hyperref}
\newcommand{\ie}{i.e.\ }
\newcommand{\cf}{cf.\ }
\newcommand{\resp}{respectively\ }
\newcommand{\Subsection}{\subsection}
\providecommand{\nobreakdash}{\nobreak\hbox{-}}
\setlength{\emergencystretch}{2em}
\multlinegap0pt
\usepackage{bm}
\usepackage{bbm}
\usepackage{dsfont}
\usepackage{xspace}
\usepackage{cancel}
\usepackage{mathtools}
\usepackage[shortlabels]{enumitem}
\usepackage{amsthm}
\makeatletter
\@ifundefined{thm}{\newtheorem{thm}{Thm}}{}
\@ifundefined{cor}{\newtheorem{cor}[thm]{Corollary}}{}
\@ifundefined{lem}{\newtheorem{lem}[thm]{Lemma}}{}
\@ifundefined{prop}{\newtheorem{prop}[thm]{Proposition}}{}
\makeatother
\title[Birkhoff center and recurrent behavior of differentially pos]{Birkhoff center and recurrent behavior of differentially positive systems on a homogeneous space}
\author{Lin Niu, Yi Wang, Yufeng Zhang}
\date{Source: arXiv:2608.14980v1 (CC BY 4.0)}
\begin{document}
\maketitle

\noindent\emph{Source and licence.} Birkhoff center and recurrent behavior of differentially positive systems on a homogeneous space. Version arXiv:2608.14980v1 (CC BY 4.0), distributed under the Creative Commons Attribution 4.0 International licence, as recorded in the arXiv licence metadata for this exact version and shown on the abstract page https://arxiv.org/abs/2608.14980v1. Full rights notice: licenses/arxiv-2608.14980-CC-BY-4.0.txt.\par\smallskip

\noindent\textit{Editorial scope.} Counted unit: the Lem whose source label is \texttt{limit\_set\_dichotomy} in the article (source0.tex lines 812--817, source label \texttt{limit\_set\_dichotomy}), delivered together with the article's own complete original proof of it (source0.tex lines 906--957, one \texttt{proof} environment of 52 lines). No mathematics is written, repaired or re-derived: statement and proof are the article's own text line for line. The proof is detached from the statement in the source: the article states the result at lines 812--817 and prints its proof at lines 906--957, and the proof environment's own header names the statement, so the association is the article's own. Required context retained uncounted: quasi-closed-prop (lines 376--378); ou (lines 605--607); sep\_n\_1 (lines 821--826); sep\_n\_2 (lines 856--861); absorb (lines 892--898). Every definition, display and label of those lines is retained unchanged. Omitted from this package and not redistributed: 1--375, 379--604, 608--811, 818--820, 827--855, 862--891, 899--905, 958--1196. Nothing inside a retained excerpt is omitted or altered: every retained body is delivered exactly as the article's lines are written, so no omission receipt of any kind accompanies this package. Citations: the retained excerpts carry no citation command at all, so no bibliography entry of the article is transcribed and no reference is redistributed. Reference adaptation: none was needed and none was made; every reference inside the delivered unit resolves inside it, so no unresolved reference or citation is produced. Printed slips kept literal: no printed word, symbol, hypothesis or number of the counted statement or of its proof was altered. Layout and class: the article's own class is \texttt{article} with packages including lineno, amsmath, amssymb, mathrsfs, amsfonts, amsthm, pstricks, pst-node, pst-text, pst-3d, psfrag, graphicx, indentfirst, enumerate; the wrapper is the lane's pinned \texttt{amsart} editorial wrapper at 10pt on A4, which restates the theorem-like environments the retained text needs and adds the article's own macro definitions that the retained text needs (none), with extra wrapper packages (bm, bbm, dsfont, xspace, cancel, mathtools, enumitem). The delivered numbering is the unit's own. Assets: no retained span contains a figure environment, a graphics-inclusion command, a PDF-page inclusion, a table or a TikZ picture; no image file is copied into this package, the package declares no asset dependency and no image file is redistributed with it. Licence: version arXiv:2608.14980v1 of this article is distributed under the Creative Commons Attribution 4.0 International licence, as recorded in the arXiv licence metadata for this exact version and shown on the abstract page https://arxiv.org/abs/2608.14980v1. A full rights notice is delivered at licenses/arxiv-2608.14980-CC-BY-4.0.txt. Characters: every retained excerpt is pure ASCII, so the delivered engine, XeLaTeX with its default Latin Modern text font, reports no missing character.
\begin{prop}\label{quasi-closed-prop}
	If the conal order ``$\leq_M$" is quasi-closed, then $\partial I^{+}(p)\subset J^{+}(p)\backslash I^{+}(p)$ and $\partial I^{-}(p)\subset J^{-}(p)\backslash I^{-}(p)$.
\end{prop}	

\begin{prop}\label{ou}
	Let $x\in \mathcal{R}(\varphi)$. Then $\omega(x)$ is either strongly ordered or unordered.
\end{prop}

	\begin{prop}\label{sep_n_1}
		Let $x\in \mathcal{R}(\varphi)$ with $x\in \partial I^{+}(x)\cup\partial I^{-}(x)$. Then, for any $y\in \mathcal{R}(\varphi)$, 
		\begin{equation*}
			\omega(y)\cap(\partial I^{+}(x)\backslash\{x\})=\emptyset \text{~and~} \omega(y)\cap(\partial I^{-}(x)\backslash\{x\})=\emptyset.
		\end{equation*}
	\end{prop}

    \begin{prop}\label{sep_n_2}
		Let $x\in \mathcal{R}(\varphi)$ with $x\in I^{+}(x)\cup I^{-}(x)$. Then for any $y\in \mathcal{R}(\varphi)$,
        \begin{equation*}
			\omega(y)\cap\partial I^{+}(x)=\emptyset \text{~and~} \omega(y)\cap\partial I^{-}(x)=\emptyset.
		\end{equation*}
	\end{prop}

	\begin{cor}\label{absorb}
		Let $x,y\in \mathcal{R}(\varphi)$ with $x\notin\omega(y)$. Whether $x\in \partial I^{+}(x)\cup \partial I^{-}(x)$ or $x\in I^{+}(x)\cup I^{-}(x)$, the following results hold: 
        \begin{enumerate}[{\rm (i)}]
            \item If $\omega(y)\cap I^{+}(x)\neq\emptyset$, then $\omega(y)\subset I^{+}(x)$;
            \item If $\omega(y)\cap I^{-}(x)\neq\emptyset$, then $\omega(y)\subset I^{-}(x)$.
        \end{enumerate}
	\end{cor}

\begin{lem}\label{limit_set_dichotomy}
	Let $x,y\in \mathcal{R}(\varphi)$. Then either \begin{enumerate}[{\rm (i)}]
	\item $\omega(x)$, $\omega(y)$ are strongly ordered, or
    \item $\omega(x)$, $\omega(y)$ are unordered.
    \end{enumerate}
\end{lem}

	\begin{proof}[Proof of Lemma \ref{limit_set_dichotomy}]
    We first consider the case that one of $x,y$ is an equilibrium. Without loss of generality, we assume $x$ is an equilibrium, then $x=\omega(x)$. If $x\in \omega(y)$, Proposition \ref{ou} implies that any points in $\omega(y)$ are ordered with $\omega(x)(=x)$. Therefore, we obtain the lemma.
    If $x\notin \omega(y)$, it follows from Propositions \ref{sep_n_1} and \ref{sep_n_2} that $\omega(y)\cap(\partial I^{+}(x)\cup\partial I^{-}(x))=\emptyset$. Then Corollary \ref{absorb} entails that one of the following alternatives must occur: {\rm (1)} $\omega(y)\subset I^{+}(x)$; {\rm (2)} $\omega(y)\subset I^{-}(x)$; {\rm (3)} $\omega(y)\subset M\backslash \overline{I^{+}(x)\cup I^{-}(x)}$.
    So, the lemma follows. In the remainder, we assume both $x,y\in \mathcal{R}(\varphi)$ are not equilibrium.

	(i) $\omega(x)\cap\omega(y)\neq\emptyset$. Suppose $a\in \omega(x)\cap\omega(y)$. Proposition \ref{ou} shows that $\omega(x)$ is either strongly ordered or unordered, as well as $\omega(y)$. Thus, if $\omega(x)\subset\omega(y)$ or $\omega(y)\subset\omega(x)$, then the lemma can be obtained directly. So in the remainder of this case, we assume $x\notin \omega(y)$ and $y\notin \omega(x)$. We consider the following two situations: (ia) $\omega(x)$ is strongly ordered; (ib) $\omega(x)$ is unordered.

    (ia) $\omega(x)$ is strongly ordered. We claim that $\omega(y)$ is strongly ordered. If the claim holds, we prove the conclusion of the lemma. Since $a\in \omega(y)$, it follows from the proof of Proposition \ref{ou} that $a\in I^{+}(y')\cup I^{-}(y')$ for any $y'\in \omega(y)\backslash\{a\}$. 
    This, with Corollary \ref{absorb}, implies that $\omega(x)\subset I^{+}(y')\cup I^{-}(y')$ for any point $y'\in \omega(y)\backslash\omega(x)$. 
    As for $y'\in \omega(y)\cap\omega(x)$, it also follows from the proof of Proposition \ref{ou} that $\omega(x)\backslash \{y'\}\subset I^{+}(y')\cup I^{-}(y')$. 
    Thus, we obtain $\omega(x)$, $\omega(y)$ are strongly ordered. Now, we prove the claim. 
    In fact, if $\omega(y)$ is unordered, then there is $y_1\in \omega(y)$ such that $a\notin J^{+}(y_1)\cup J^{-}(y_1)$. So, $a\notin \overline{I^{+}(y_1)}\cup \overline{I^{-}(y_1)}$.
    Choose a point $x_1\in \omega(x)$ which is close enough to $a$ that $x_1\notin\overline{I^{+}(y_1)}\cup \overline{I^{-}(y_1)}$. 
    Together with the proof of Proposition \ref{ou}, the fact that $x_1,a\in \omega(x)$ implies that $x_1\in I^{+}(a)\cup I^{-}(a)$. 
    Thus, there is a point $x_{\tau_1}$ in $O(x)$ close enough to $x_1$ such that $x_{\tau_1}\in I^{+}(a)\cup I^{-}(a)$ and $x_{\tau_1}\notin\overline{I^{+}(y_1)}\cup \overline{I^{-}(y_1)}$. Note the fact that $x_{\tau_1}\in I^{+}(a)$ ($x_{\tau_1}\in I^{-}(a)$, respectively) implies $a\in I^{-}(x_{\tau_1})$ ($a\in I^{+}(x_{\tau_1})$, respectively) which implies $\omega(y)\cap (I^{+}(x_{\tau_1})\cup I^{-}(x_{\tau_1})) \neq\emptyset$ as $a\in \omega(y)$. Since $x_{\tau_1}\notin\omega(y)$ (as $x_{\tau_1}\in O(x)$ and $x\notin\omega(y)$), it has $\omega(y)\subset I^{+}(x_{\tau_1})\cup I^{-}(x_{\tau_1})$ by Corollary \ref{absorb}. 
    Then as $y_1\in \omega(y)$, we have $y_1\in I^{+}(x_{\tau_1})\cup I^{-}(x_{\tau_1})$, and hence, $x_{\tau_1}\in I^{+}(y_1)\cup I^{-}(y_1)$, which is a contradiction to $x_{\tau_1}\notin\overline{I^{+}(y_1)}\cup \overline{I^{-}(y_1)}$. So, we obtain the claim.
    
    
    (ib) $\omega(x)$ is unordered. Then $\omega(y)$ is unordered. (We prove it by contradiction. If $\omega(y)$ is strongly ordered, then there is $y_2\in \omega(y)$ such that $a\in I^{+}(y_2)\cup I^{-}(y_2)$. Recall that $a\in \omega(x)\cap\omega(y)$, one can choose a point $x_2\in \omega(x)$ close enough to $a$ that $x_2\in I^{+}(y_2)\cup I^{-}(y_2)$. Since $\omega(x)$ is unordered, it also has $x_2\notin\overline{I^{+}(a)}\cup \overline{I^{-}(a)}$. Thus, there is a point $x_{\tau_2}$ in $O(x)$ which is close enough to $x_2$ such that $x_{\tau_2}\in I^{+}(y_2)\cup I^{-}(y_2)$ and $x_{\tau_2}\notin\overline{I^{+}(a)}\cup \overline{I^{-}(a)}$. Then by a similar argument in situation (ia) as above, one can obtain a contradiction.) Since $a\in \omega(x)\cap\omega(y)$, it has $a\notin \overline{I^{+}(y')}\cup \overline{I^{-}(y')}$ for any $y'\in \omega(y)\backslash\{a\}$. This, together with Corollary \ref{absorb}, implies that $\omega(x) \cap (I^{+}(y')\cup I^{-}(y'))=\emptyset$ for any point $y'\in \omega(y)\backslash \omega(x)$. As for $y'\in \omega(y)\cap \omega(x)$, the fact that $\omega(x)$ is unordered implies that $(\omega(x)\backslash \{y'\}) \cap (I^{+}(y')\cup I^{-}(y'))=\emptyset$. So, we obtain that $(\omega(x)\backslash \{y'\}) \cap (I^{+}(y')\cup I^{-}(y'))=\emptyset$ for all $y'\in \omega(y)$. Thus, $\omega(x)$, $\omega(y)$ are unordered. (In fact, if there exist $p\in \omega(x)$ and $q\in \omega(y)$ such that $p<_M q$, then the strong positivity of the system entails that $\varphi_{t}(p)\in I^{-}(\varphi_{t}(q))$ for $t>0$, a contradiction.)

    (ii) $\omega(x)\cap\omega(y)=\emptyset$. We assert that \emph{there is no points $x_{0}\in\omega(x)$ and $y_{0}\in\omega(y)$ such that $y_{0}\in \partial I^{+}(x_{0})\cup \partial I^{-}(x_{0})$}. 
	Before proving this assertion, we show how it implies the conclusion of this lemma.
	In fact, choose any $x_*\in\omega(x)$, $y_*\in\omega(y)$. Since $\omega(x)\cap\omega(y)=\emptyset$, one has $x_*\neq y_*$. We consider the following two case:
	(iia) $y_*\in I^{+}(x_*)\cup I^{-}(x_*)$; (iib) $y_1\notin I^{+}(x_*)\cup I^{-}(x_*)$.
		
	(iia) $y_*\in I^{+}(x_*)\cup I^{-}(x_*)$. Then we have
	\begin{equation}\label{x1_approx_omegay}
		\omega(y)\subset I^{+}(x_*)\cup I^{-}(x_*).
	\end{equation}
	(Otherwise, one can find $y_3\in\omega(y)$ such that $y_3\notin I^{+}(x_*)\cup I^{-}(x_*)$. Noting $y_*\in I^{+}(x_*)\cup I^{-}(x_*)$, it follows from the connectedness of $\omega(y)$ that there is $y_4\in\omega(y)$ such that $y_4\in\partial I^{+}(x_*)\cup \partial I^{-}(x_*)$, contradicting the assertion.)
	Consequently, \eqref{x1_approx_omegay} implies that $\omega(y)\subset I^{+}(x')\cup I^{-}(x')$ for any $x'\in \omega(x)$.
	(Otherwise, there exist $x_3\in\omega(x)$ and $y_4\in\omega(y)$ such that $y_4\notin I^{+}(x_3)\cup I^{-}(x_3)$, equivalently $x_3\notin I^{+}(y_4)\cup I^{-}(y_4)$. But, note $x_*\in I^{+}(y_4)\cup I^{-}(y_4)$ by \eqref{x1_approx_omegay}, the connectedness of $\omega(x)$ implies that there is a $x_4\in\omega(x)$ such that $x_4\in \partial I^{+}(y_4)\cup \partial I^{-}(y_4)$, which contradicts the assertion). Thus, it has $\omega(x)$, $\omega(y)$ are strongly ordered.
		
	(iib) $y_*\notin I^{+}(x_*)\cup I^{-}(x_*)$. One can obtain that 
    \begin{equation}\label{x2_approx_omegay}
		\omega(y)\cap (I^{+}(x_*)\cup I^{-}(x_*))=\emptyset.
	\end{equation}
    (Otherwise, one can find $y_5\in\omega(y)$ such that $y_5\in I^{+}(x_*)\cup I^{-}(x_*)$. Note also that $y_*\notin I^{+}(x_*)\cup I^{-}(x_*)$. Then it follows from the connectedness of $\omega(y)$ that there is $y_6\in\omega(y)$ such that $y_6\in \partial I^{+}(x_*)\cup \partial I^{-}(x_*)$, contradicting the assertion.) 
    Then, for any $x'\in\omega(x)$, it has 
    \begin{equation}\label{x3_approx_omegay}
        \omega(y)\cap (I^{+}(x')\cup I^{-}(x'))=\emptyset.
    \end{equation} (Otherwise, there exist $x_5\in\omega(x)$ and $y_6\in\omega(y)$ such that $y_6\in I^{+}(x_5)\cup I^{-}(x_5)$, that is, $x_5\in I^{+}(y_6)\cup I^{-}(y_6)$. Note also $y_6\notin I^{+}(x_*)\cup I^{-}(x_*)$ by \eqref{x2_approx_omegay}, that is, $x_*\notin I^{+}(y_6)\cup I^{-}(y_6)$. 
    Then, by the connectedness of $\omega(x)$, one can obtain $x_6\in\omega(x)$ such that $x_6\in \partial I^{+}(y_6)\cup \partial I^{-}(y_6)$, a contradiction to the assertion). Furthermore, we show that $\omega(y)\cap (\overline{I^{+}(x')}\cup \overline{I^{-}(x')})=\emptyset$ for any $x'\in\omega(x)$. (If not, there exist $x_7\in\omega(x)$ and $y_7\in\omega(y)$ such that $y_7\in \overline{I^{+}(x_7)}\cup \overline{I^{-}(x_7)}$ with $x_7\neq y_7$, and then $\varphi_t(y_7)\in I^{+}(\varphi_t(x_7))\cup I^{-}(\varphi_t(x_7))$ for any $t>0$ by (H2) (or rather, Proposition \ref{quasi-closed-prop}) and the strong positivity of the system, a contradiction to \eqref{x3_approx_omegay}). So, $\omega(x)$, $\omega(y)$ are unordered.
    
    Therefore, we have proved the lemma provided that the assertion is confirmed.
		
	Finally, it remains to prove the assertion in (ii). Suppose that there exist $x_{0}\in\omega(x)$ and $y_{0}\in\omega(y)$ such that $y_{0}\in \partial I^{+}(x_{0})\cup \partial I^{-}(x_{0})$ with $x_0\neq y_0$. 
    Then by (H2) (or rather, Proposition \ref{quasi-closed-prop}) and the strong positivity of the system, we obtain that $\varphi_1(y_{0})\in I^{+}(\varphi_1(x_{0}))\cup I^{-}(\varphi_1(x_{0}))$. Since $\varphi_1(x_{0})\in \omega(x)$, there are $x_8\in O^+(x)$ close enough to $\varphi_1(x_{0})$ that $\varphi_1(y_{0})\in I^{+}(x_8)\cup I^{-}(x_8)$ by the continuity of $I^{+}$ and $I^{-}$. 
    Also since $\varphi_1(y_{0})\in \omega(y)$, choose $y_8\in O^+(y)$ close enough to $\varphi_1(y_{0})$ that $y_8\in I^{+}(x_8)\cup I^{-}(x_8)$. Since $x_8\notin \omega(y)$ (as $\omega(x)\cap\omega(y)=\emptyset$), according to Corollary \ref{absorb}, it has $\omega(y)\subset I^{+}(x_8)\cup I^{-}(x_8)$. 
    Thus, for any $y'\in \omega(y)$, it has  $y'\in I^{+}(x_8)\cup I^{-}(x_8)$, that is, $x_8\in I^{+}(y')\cup I^{-}(y')$. 
    So, it has $\omega(x)\subset I^{+}(y')\cup I^{-}(y')$ for any $y'\in \omega(y)$ by Corollary \ref{absorb} and the fact $y'\notin\omega(x)$ (as $\omega(x)\cap\omega(y)=\emptyset$). And hence, $x_0\in I^{+}(y_0)\cup I^{-}(y_0)$, that is $y_0\in I^{+}(x_0)\cup I^{-}(x_0)$, which contradicts with $y_{0}\in\partial I^{+}(x_{0})\cup \partial I^{-}(x_{0})$. Thus, we have proved the assertion and completed the proof.
\end{proof}

\end{document}
